\fsection{Ej 1 Pág 110}
\fsubsection
[\textcolor{waveBlue2}{INVESTIGACIÓN.}]
{\textcolor{waveBlue2}{\textbf{\textit{INVESTIGACIÓN.}}}}
{
	\textbf{\textit{Tarea.}}
	Business intelligence, o inteligencia empresarial, comúnmente conocido como BI, es un conjunto de estrategias,
	aplicaciones, etc., cuyo objetivo es administrar y crear conocimiento a partir del análisis de los datos de una
	empresa.
}\\

Actualmente, el BI es esencial tanto en la Administración pública como en el sector productivo en general,
y vamos a comprobarlo mediante una investigación de tendencias.


\begin{solution}
	\textbf{Tendencia elegida: Inteligencia Artificial y \textit{Machine Learning} aplicados al BI}

	\begin{itemize}
		\item \boldit{Definición:} Integración de algoritmos de IA y ML en plataformas de BI para
		      automatizar el análisis de datos, detectar patrones ocultos y generar predicciones sin
		      intervención manual.

		\item \boldit{Aplicaciones prácticas:}
		      \begin{itemize}
			      \item Detección automática de anomalías en datos financieros u operativos.
			      \item Segmentación inteligente de clientes.
			      \item Previsión de demanda y ventas.
			      \item Generación automática de informes y dashboards.
		      \end{itemize}

		\item \boldit{Aplicación al sector:}
		      \begin{itemize}
			      \item Predicción de fallos en infraestructura y servidores.
			      \item Monitorización inteligente de redes.
			      \item Detección de intrusiones y amenazas de seguridad mediante ML.
			      \item Optimización automática de recursos en entornos cloud.
		      \end{itemize}

		\item \boldit{Beneficios:}
		      \begin{itemize}
			      \item Decisiones más rápidas y basadas en datos.
			      \item Reducción de errores humanos en el análisis.
			      \item Capacidad de procesar grandes volúmenes de datos en tiempo real.
			      \item Automatización de tareas repetitivas de análisis.
		      \end{itemize}

		\item \boldit{Desafíos:}
		      \begin{itemize}
			      \item Necesidad de datos de calidad y en cantidad suficiente.
			      \item Riesgo de sesgos en los modelos.
			      \item Requiere personal con formación especializada.
			      \item Coste elevado de implementación inicial.
		      \end{itemize}

		      \begin{itemize}
			      \item \boldit{Empresas:}
			            \begin{itemize}
				            \item \textbf{Microsoft} --- Power BI con Azure Machine Learning.
				            \item \textbf{Google} --- Looker integrado con Vertex AI.
				            \item \textbf{Salesforce} --- Einstein Analytics con ML integrado.
				            \item \textbf{IBM} --- Watson Analytics para BI empresarial.
			            \end{itemize}
		      \end{itemize}
	\end{itemize}

\end{solution}
